\section{The endomorphism algebra of the residual complex}\label{sec:I}
The finite topology on $\Hom_A(M,N)$ means pointwise convergence
when the modules are discrete. The universal problem below uses
separated complete graded topological $A$-algebras, continuous
multiplication, and continuous graded homomorphisms, with $A$ discrete
as base. The prescribed degree zero has the product of local adic
topologies. This specifies an interpretation of completion; the
uncertain transcription on 43:65 does not uniquely specify a category.

\begin{proposition}[I1: degree-zero diagonality; Proved]\label{i:zero}
For a regular noetherian local ring, normalize the residual complex by
$K^i=\bigoplus_{\operatorname{ht}\mathfrak p=i}E_A(A/\mathfrak p)$.
Then $\End_A^0(K)=\prod_{\mathfrak p}\widehat{A_{\mathfrak p}}$
as topological algebras, and $\End_A^i(K)=0$ for $i<0$.
\end{proposition}
\begin{proof}
Write $E_{\mathfrak p}=E_A(A/\mathfrak p)$. If
$a\in\mathfrak p\setminus\mathfrak q$, every element of
$E_{\mathfrak p}$ is killed by a power of $a$, while multiplication
by $a$ is invertible on $E_{\mathfrak q}$. Thus the corresponding Hom
vanishes unless $\mathfrak p\subseteq\mathfrak q$. Distinct primes
of the same height are incomparable; projections detect every element
of a direct sum. This proves diagonality and negative-degree vanishing.
Matlis' endomorphism theorem gives the displayed rings
\cite[Lemma 2, proof, p.~123]{Sastry1992}. Annihilators of finite
subsets of $E_{\mathfrak p}$ give exactly the local adic topology:
the submodules killed by powers of the maximal ideal have finite
length over $A_{\mathfrak p}$ and exhaust the hull. The direct sum
gives the product topology.
\end{proof}

\begin{theorem}[I2: complete surface presentation; Proved]\label{i:surface}
Let $A=k[[s,t]]$, $F=\operatorname{Frac}(A)$, and
$\mathfrak m=(s,t)$. Then $\End_A^*(K)$ with differential $\theta$
is universal among the separated complete graded topological
$A$-algebras with degree zero
\[
 B^0=F\times\prod_{\operatorname{ht}\mathfrak p=1}
                 \widehat{A_{\mathfrak p}}\times A,
\]
one specified element of degree one, its square zero, and the
immediate-specialization incidence relations. In particular the
canonical map from the algebra defined by this precise universal
property is an isomorphism, including degree two.
\end{theorem}
\begin{proof}
Let $\mathcal P_1$ denote the height-one primes, put
$R_{\mathfrak p}=\widehat{A_{\mathfrak p}}$, choose a generator
$f_{\mathfrak p}$ of $\mathfrak p$, and put
$F_{\mathfrak p}=\operatorname{Frac}(R_{\mathfrak p})$.
The factorial Gorenstein ring $A$ has residual resolution
\begin{equation}\label{i:res}
 0\to A\to F\xrightarrow{d^0}
 I=\bigoplus_{\mathfrak p\in\mathcal P_1}E_{\mathfrak p}
 \xrightarrow{d^1}E\to0,
\end{equation}
where $E=E_A(k)$, $E_{\mathfrak p}=F/A_{\mathfrak p}$, and
$d^0$ consists of quotient maps. Write $\delta_{\mathfrak p}$ for
the components of $d^1$ and $D(-)=\Hom_A(-,E)$.

\emph{Individual arrows.}
By \cite[Theorem 1.1, p.~1]{Schenzel2013}, $D(E_{\mathfrak p})$
is free of rank one over $R_{\mathfrak p}$: $A/\mathfrak p$ is
complete of dimension one. Its element $\delta_{\mathfrak p}$ is
a basis. Indeed, applying $\Hom_A(A/\mathfrak p,-)$ to
\eqref{i:res} leaves only
$\kappa(\mathfrak p)\to\Hom_A(A/\mathfrak p,E)$.
The free resolution of $A/(f_{\mathfrak p})$ shows this map is
surjective with kernel isomorphic to $A/\mathfrak p$, so it is
nonzero. Dualizing
$0\to\kappa(\mathfrak p)\to E_{\mathfrak p}
\xrightarrow{f_{\mathfrak p}}E_{\mathfrak p}\to0$
shows that the kernel of restriction of $D(E_{\mathfrak p})$ to
$D(\kappa(\mathfrak p))$ is
$f_{\mathfrak p}D(E_{\mathfrak p})$. Hence
$\delta_{\mathfrak p}$ is a unit times a basis. We use it to identify
\begin{equation}\label{i:arrows}
 \Hom_A(E_{\mathfrak p},E)=R_{\mathfrak p},\qquad
 \Hom_A(F,E_{\mathfrak p})=F_{\mathfrak p}.
\end{equation}
For the second identity, apply $\Hom_{A_{\mathfrak p}}(-,E_{\mathfrak p})$
to $0\to A_{\mathfrak p}\to F\to E_{\mathfrak p}\to0$
and localize at $f_{\mathfrak p}$. The middle Hom is already an
$F$-vector space. Explicitly $a\in F_{\mathfrak p}$ sends
$f\in F$ to $fa\bmod R_{\mathfrak p}$, using
$E_{\mathfrak p}=F_{\mathfrak p}/R_{\mathfrak p}$.

\emph{All graded pieces.}
Put
\[
 P=\prod_{\mathfrak p}R_{\mathfrak p},\qquad
 \mathbb A=\left\{(a_{\mathfrak p})\in\prod_{\mathfrak p} F_{\mathfrak p}:
 a_{\mathfrak p}\in R_{\mathfrak p}\text{ for almost all }\mathfrak p\right\}.
\]
For $S=A\setminus\{0\}$, factoriality gives $\mathbb A=S^{-1}P$:
clear its finitely many negative valuations by a product of powers of
$f_{\mathfrak p}$. Conversely every denominator has finitely many
prime divisors. A family $F\to E_{\mathfrak p}$ has values in the
direct sum precisely when its coefficients are integral almost
everywhere: evaluate at $1$ for necessity and use the finitely many
poles of each $f\in F$ for sufficiency. Consequently
\begin{equation}\label{i:pieces}
 \End_A^1(K)=\mathbb A\oplus P,\qquad
 \End_A^2(K)=\mathbb A/F.
\end{equation}
For the second equality dualize \eqref{i:res}, using $D(E)=A$:
\begin{equation}\label{i:dual}
 0\to A\to P\to D(F)\to E\to0.
\end{equation}
The first map is diagonal, since its components are
$a\delta_{\mathfrak p}$. Localizing kills $E$ and leaves $D(F)$
unchanged, proving the equality. For $c\in P$ representing $I\to E$
and $a\in\mathbb A$ representing $F\to I$, composition is
\begin{equation}\label{i:product}
 c\circ a=[(c_{\mathfrak p}a_{\mathfrak p})_{\mathfrak p}]
                         \in\mathbb A/F.
\end{equation}
All other products of positive-degree arrows vanish by endpoints;
degrees above two vanish. The differential is $(1,1)$, with square
$[1]=0$. Thus the degree-two relation is precisely diagonal $F$,
not an independent zero relation for each intermediate prime.

\emph{The topologies in degree one.}
The first summand has the usual restricted-product topology: $P$ with
its product adic topology is open. Indeed, vanishing at $1/u$ means
membership in $uP$, and these form a basis. The second summand has
the product of topologies $\tau_{\mathfrak p}$ with bases
\begin{equation}\label{i:mixed}
 uA+f_{\mathfrak p}^n R_{\mathfrak p}\ (u\notin\mathfrak p),
 \quad\text{or equivalently}\quad
 \mathfrak m^N A+f_{\mathfrak p}^n R_{\mathfrak p}.
\end{equation}
To prove this, exactness and factoriality give
\[
 \ker\delta_{\mathfrak p}
  =(A[1/f_{\mathfrak p}]+A_{\mathfrak p})/A_{\mathfrak p}.
\]
A fraction whose other components vanish has no other prime factors
in its denominator. Any finite set in $E_{\mathfrak p}$ lies in the
cyclic $A$-module generated by $1/(u f_{\mathfrak p}^n)$ for some
$u\notin\mathfrak p$. Its annihilator in the Hom space is exactly
$uA+f_{\mathfrak p}^nR_{\mathfrak p}$: modulo
$f_{\mathfrak p}^n$, the element $c/u$ must come from $A$.
Since $(u,f_{\mathfrak p}^n)$ is $\mathfrak m$-primary, some
$\mathfrak m^N$ is contained in it; conversely choose
$u\in\mathfrak m^N\setminus\mathfrak p$. This proves the
equivalent basis. At $\mathfrak p=(s)$, it is coefficientwise
$t$-adic convergence on $k((t))[[s]]$. Hence
$t^N\delta_{\mathfrak p}\to0$, although $t^N$ is always an
$s$-adic unit. The two occurrences of $P$ require different topologies.

\emph{Degree two.}
Write $P_\tau$ for the product with \eqref{i:mixed}. The subgroup
$H=\{h\in D(F):h(1)=0\}$ is open and topologically $P_\tau/A$.
For the quotient assertion, if $M'\subset M$ and $E$ is injective,
restriction $D(M)\to D(M')$ is an open surjection for the finite
topologies over a noetherian ring. The image of the annihilator of a
finitely generated $N\subset M$ is exactly the annihilator of
$M'\cap N$: extend a map from $M'$ by zero on $N$, and then use
injectivity. The intersection is finitely generated, proving openness.
Apply this to $F/A\subset I$; the kernel is $D(E)=A$.
Furthermore $D(F)=\bigcup_{u\ne0}u^{-1}H$, since every element of
$E$ is killed by a nonzero element of $A$, and multiplication by $u$
is a topological automorphism. This specifies the topology on
$\mathbb A/F$ without an unjustified ordinary-adelic quotient.

\emph{The universal morphism.}
Let $C$ have the prescribed degree zero and $T\in C^1$ satisfying
the relations. Write $e_\eta,e_P,e_m$ for the three block
idempotents. Finite sums of the individual height-one idempotents
converge to $e_P$, so continuity and incidence give
\begin{equation}\label{i:T}
 T=e_PTe_\eta+e_mTe_P.
\end{equation}
For $a=b/u\in\mathbb A$, $b\in P$, $u\ne0$, and $c\in P$, set
\begin{align}\label{i:normal}
 \alpha_C(a)&=bTe_\eta u_\eta^{-1},&
 \beta_C(c)&=e_mTc,\\
 \gamma_C([a])&=e_mTbTe_\eta u_\eta^{-1}.\nonumber
\end{align}
The elements of $P$ occupy the height-one block, and $u_\eta^{-1}$
means $1/u$ in the generic block. Clearing denominators and using
$A$-linearity proves independence of representations. For diagonal
$a=v/u\in F$, the last expression is
$v e_mTe_PTe_\eta u_\eta^{-1}=0$, by \eqref{i:T} and $T^2=0$.
Thus $\gamma_C$ descends to $\mathbb A/F$.
These formulas give an algebra homomorphism, the identity in degree
zero: the only composable positive-degree product is
$\beta_C(c)\alpha_C(a)=\gamma_C([ca])$, agreeing with
\eqref{i:product}. All others vanish by endpoints. The diagonal
actions agree directly, using $A$-linearity for the closed-point
action and the right $F$-action at the generic point.

We check continuity carefully. The map $\alpha_C$ is continuous
on the open subgroup $P$ with product adic topology, and hence on all
$u^{-1}P$. For $\beta_C$, given a zero-neighborhood $U$, choose
$V$ with $V+V\subset U$. Continuity from product adic degree zero
gives a finite set $T_0$ of primes and powers
$f_{\mathfrak p}^{n_{\mathfrak p}}$ such that tuples satisfying
those divisibilities map into $V$; all remaining coordinates are
unrestricted. Each map
$A\to C^1$, $a\mapsto a e_mTe_{\mathfrak p}$, is continuous
for the $\mathfrak m$-adic topology. Choose $N$ so that the sum of
its images of $\mathfrak m^N$ over $\mathfrak p\in T_0$ lies
in $V$. Tuples satisfying
$c_{\mathfrak p}\in\mathfrak m^N A+
f_{\mathfrak p}^{n_{\mathfrak p}}R_{\mathfrak p}$ split into a
tuple controlled by the first condition and finitely many such
$\mathfrak m^N A$ terms. Their images lie in $U$. This proves
continuity from $P_\tau$.
The continuous map $c\mapsto\beta_C(c)Te_\eta$ vanishes on
diagonal $A$, hence descends to the open quotient $H=P_\tau/A$.
Multiplication by $u_\eta^{-1}$ gives continuity of $\gamma_C$
on every $u^{-1}H$ and globally.

The morphism sends $\theta$ to $T$. It is unique because
\eqref{i:normal}, with $T=\theta$, represents every positive-degree
endomorphism. Finally, all Hom spaces in the finite topology are
closed subgroups of products of discrete target modules, defined by
closed $A$-linearity equations, so are separated and complete.
Composition is jointly continuous by controlling finitely many inputs
and their finitely many intermediate images. Statement~\ref{i:zero}
identifies the degree-zero topology. This proves the theorem.
\end{proof}

\paragraph{Remaining gate.}
This proves the requested complete surface case, including degree two,
for the stated completion convention. It does not yet prove the
presentation for every essentially finite type regular local ring or
in every dimension. In a noncomplete surface the rank-one Matlis-dual
input can fail: Schenzel's Theorem 1.1 makes completeness of
$A/\mathfrak p$ its exact condition. Longer residual resolutions
require a proof that all higher relations follow from $\theta^2$.
Constructing a residual complex or proving its residue theorem does
not establish that assertion. Yekutieli explicitly specifies
semi-topological rings and a tensor topology
\cite[Definitions 1.2.1 and 1.2.11, pp.~15 and 18]{Yekutieli1992};
his construction is not an opened presentation theorem for the whole
endomorphism algebra.

